% =====================================================================
%  Lugworm: Beginner's Guide
%  Plain-language introduction, no prior Grasshopper/clay-printing
%  knowledge assumed.
%
%  Compile with pdflatex (run twice for the table of contents / links):
%    pdflatex beginner-guide.tex
%    pdflatex beginner-guide.tex
% =====================================================================
\documentclass[11pt]{article}

\usepackage[a4paper,margin=2.2cm]{geometry}
\usepackage[utf8]{inputenc}
\usepackage[T1]{fontenc}
\usepackage{textcomp}
\usepackage{amsmath}
\usepackage{booktabs}
\usepackage{longtable}
\usepackage{array}
\usepackage{xcolor}
\usepackage{enumitem}
\usepackage{parskip}
\usepackage{titlesec}
\usepackage{fancyhdr}
\usepackage{hyperref}

\definecolor{lugwormblue}{RGB}{0,70,140}
\hypersetup{
  colorlinks=true,
  linkcolor=lugwormblue,
  citecolor=lugwormblue,
  urlcolor=lugwormblue,
  pdftitle={Lugworm Beginner's Guide},
  pdfauthor={Jasper Vollmert, IDAC, Hamburg University of Technology},
  bookmarksopen=true
}

\titleformat{\section}{\normalfont\Large\bfseries\color{lugwormblue}}{\thesection}{0.6em}{}
\titleformat{\subsection}{\normalfont\large\bfseries}{\thesubsection}{0.6em}{}
\titleformat{\subsubsection}{\normalfont\normalsize\bfseries}{\thesubsubsection}{0.6em}{}

\pagestyle{fancy}
\fancyhf{}
\fancyhead[L]{\small Lugworm: Beginner's Guide}
\fancyhead[R]{\small\leftmark}
\fancyfoot[C]{\thepage}
\renewcommand{\headrulewidth}{0.4pt}

\newcommand{\param}[1]{\texttt{#1}}
\newcommand{\iamdeg}{$^{\circ}$}

\setlist[itemize]{leftmargin=1.4em, itemsep=0.15em, topsep=0.2em}
\setlist[enumerate]{leftmargin=1.6em, itemsep=0.15em, topsep=0.2em}

\title{\color{lugwormblue}\textbf{Lugworm}\\[0.3em]\Large Beginner's Guide}
\author{Jasper Vollmert \\ \small Institute of Digital and Autonomous Construction (IDAC) \\ \small Hamburg University of Technology \\ \small jasper.vollmert@tuhh.de}
\date{\small \today}

\begin{document}
\maketitle

\begin{abstract}
\noindent This guide is written for a first session with Lugworm.
No previous experience with Grasshopper or clay printing is assumed.
Lugworm is a Grasshopper plugin for slicing and printing clay and
paste materials (liquid deposition modeling, LDM) with a ram or piston
extruder, a robot arm, or both, and includes a calibration workflow.

The reference (\emph{reference.tex}) explains the reasoning behind each
setting and serves as a lookup once the basics are familiar. The
exercises (\emph{exercises.tex}) provide practice.
\end{abstract}

\tableofcontents
\newpage

% =====================================================================
\section{What the plugin does}
% =====================================================================

All components are located on one Grasshopper tab, \textbf{Lugworm}, in
numbered panels that follow the order in which the components are used:

\begin{center}
\texttt{1 Slice\quad|\quad 2 Place\quad|\quad 3 G-code\quad|\quad 4 Robot\quad|\quad 5 Analyze\quad|\quad 6 Calibrate\quad|\quad 7 Print}
\end{center}

The core pipeline for a normal print consists of three components:

\begin{verbatim}
your 3D geometry
      |
      v
+-------------+  contours (one    +-------------+  moved contours   +--------------+
|Slice |- branch/layer) ->|Place |------------------>|Lugworm G-code |--> .gcode
+-------------+                   +-------------+                   |   Generator  |
 "cut it into                      "put it on                       +--------------+
  layers"                           the bed,                         "write the machine
                                     fix the seam"                    instructions"
\end{verbatim}

The pipeline runs without a robot selected, so the plugin can be used
as an ordinary slicer for any Marlin or Duet printer.

When a robot arm does the printing, the machine is selected explicitly
and carried on one wire. \param{Robot Definition} (panel
\textbf{4 Robot}) picks the machine profile: the lab SCARA (Potterbot
Scara v4), a generic Cartesian machine, or \param{none}. Every
robot-aware component reads the profile from a single \param{robot}
input. \param{Robot Simulation} visualizes and validates the
toolpath on the selected machine.

Three further groups of components exist alongside the pipeline:

\begin{itemize}
  \item \textbf{Analyze} (\param{Overhang}, \param{Print Stats}) checks a slice or a finished G-code file before a print is committed.
  \item \textbf{Calibrate} (\param{DoE Design}, \param{RSM Evaluate}, \param{Desirability}, \param{Test Geometry}) finds settings for a new material with designed experiments rather than trial and error. The workflow is an advanced topic and is only summarised here; the reference and the exercises describe the full procedure.
  \item \textbf{Print} (the \param{Lugworm Duet} components, panel \textbf{7 Print}) controls a Duet-based machine directly over WiFi from Grasshopper: connect, jog, upload, print, and monitor, behind fail-closed safety gates. The Print group is not needed for slicing and is documented in \emph{DUET\_CONTROL.md} in \texttt{docs/reference/}.
\end{itemize}

One utility completes the tab. \param{Graph Mapper}
(\textbf{1 Slice}) remaps a number stream through an editable curve,
which is useful for flow or speed gradients.

% =====================================================================
\section{Before starting}
% =====================================================================

\begin{itemize}
  \item Rhino 8 with Grasshopper, and Lugworm installed as a Yak package (package name \param{Lugworm}, author \param{Lugworm}). Components are dragged onto the canvas from the \textbf{Lugworm} tab. To check which version is installed, open Rhino's \param{PackageManager} (the Lugworm entry shows it) or the Details tab of the \param{.gha} file's properties.
  \item \param{Slice}, \param{G-code Generator}, \param{G-code Restart}, \param{Overhang}, \param{RSM Evaluate}, and \param{Desirability} are Heavy components. Each has a manual compute button instead of recomputing on every change (see \S 6). This behaviour is intended.
  \item The Rhino document may use millimetres, inches, or metres. The plugin converts geometry at the boundary. Numeric inputs documented in mm (\param{layerHeight}, \param{width}, \param{spacing}, and so on) stay in millimetres regardless of the document units, and the G-code is always written in millimetres.
  \item Some features depend on the printer or firmware, or on data that has not been measured yet, such as the pump calibration table. The plugin reports such cases in the \param{info} outputs rather than guessing. Read the \param{info} output of every component.
\end{itemize}

% =====================================================================
\section{A first slice with a few inputs}
% =====================================================================

On \param{Slice}:

\begin{longtable}{@{}p{3.0cm}p{9.5cm}@{}}
\toprule
\textbf{Input} & \textbf{Value} \\
\midrule
\endhead
\param{geo} & The 3D shape: a mesh, Brep, or surface. A list works as well. \\
\param{layerHeight} & Start with the default of \param{3.0}\,mm. \\
\param{wallCount} & \param{1} to start. \\
\param{spacing} & The default of \param{10.0}\,mm (must be greater than 0). \\
compute button & Click once. Slice is a Heavy component. \\
\bottomrule
\end{longtable}

On \param{Place}, connect \param{contours}. The default
\param{printCenter} of (339, $-193$, 0) is the usual spot on the lab
bed, and \param{seamMode = Align} is a reasonable start. On another
machine, set \param{printCenter} directly or connect an \param{Lugworm
Print Bed} definition (see \S 4).

On \param{G-code Generator}:

\begin{longtable}{@{}p{3.0cm}p{9.5cm}@{}}
\toprule
\textbf{Input} & \textbf{Value} \\
\midrule
\endhead
\param{crvs} & Connect \param{moved} from \param{Place}. \\
\param{layerHeight} & The same value as on \param{Slice}. \\
\param{width} & The default of \param{10.0}\,mm. \\
\param{ramDiameter} & The inside diameter of the cartridge or barrel in mm (default \param{94}). \\
\param{printSpeed} & The default of \param{43}\,mm/s is a reasonable start. \\
\param{robot} & On a Cartesian printer, connect an \param{Robot Definition} with the profile \param{generic-cartesian} or \param{none}. Without a definition, the output keeps the legacy coordinate frame of the lab SCARA; see \S 4. \\
\param{savePath} & A folder for the \param{.gcode} file, or use \param{Save G-code} instead (see below). \\
compute button & Click once. The G-code Generator is a Heavy component as well. \\
\bottomrule
\end{longtable}

The \param{gcode} output holds the finished file. The \param{info}
output of every component summarises what happened in one line. Read
\S 5 before sending any file to a machine.

% =====================================================================
\section{The pipeline in detail}
% =====================================================================

\subsection{Slice: turns the shape into layers}

\begin{longtable}{@{}p{2.6cm}p{3.6cm}p{1.6cm}p{4.6cm}@{}}
\toprule
\textbf{Input} & \textbf{Meaning} & \textbf{Default} & \textbf{When to change} \\
\midrule
\endhead
\param{geo} & The shape or shapes to slice: mesh, Brep, or surface, single or list. & required & Always. \\
\param{layerHeight} & Layer thickness in mm. & 3.0 & Thicker layers print faster and are more visible. \\
\param{firstLayer} & Height of the first layer only; \param{0} means half a normal layer. & 3.0 & Raise for more first-layer squish. \\
\param{wallCount} & Number of outline passes per layer. & 1 & Raise for a thicker shell. \\
\param{spacing} & Distance between wall and infill layers in mm; must be greater than 0. & 10.0 & Match to the layer width. \\
\param{mode} & \param{0} planar, \param{1} conical (non-planar). & 0 & Leave at \param{0} until \S 7 has been read; mode 1 is advanced. \\
\param{infill} & \param{0} none, \param{1} zigzag, \param{2} concentric, \param{3} spiral, \param{4} fermat. & 0 & \param{4} (fermat) gives the smoothest and most continuous fill. \\
\param{cornerRadius} & Rounds sharp corners slightly to avoid over-fill blobs. & 0.0 & Raise a little if corners look overstuffed. \\
\param{tol} & Curve simplification tolerance in mm. & 0.05 & Rarely changed. \\
\bottomrule
\end{longtable}

Toggle groups can be added when needed: \textbf{Adaptive} (varies the
layer height with the slope), \textbf{Conical}, \textbf{Zigzag}, and
\textbf{Spiralize} (advanced, see \S 7), \textbf{Meshing} (controls how
Breps are converted to a mesh internally), and \textbf{Multi-object}
(\param{printOrder}, the sequence when a list of shapes is supplied).

Outputs: \param{contours} (a curve tree with one branch per layer, or
\param{\{object;layer\}} for a list of shapes), \param{layers} (the
count), and \param{info} (diagnostics per object, reported without
failing the whole slice).

\subsection{Place: positions the part and sets the seam}

\begin{longtable}{@{}p{2.8cm}p{5.2cm}p{4.2cm}@{}}
\toprule
\textbf{Input} & \textbf{Meaning} & \textbf{Default} \\
\midrule
\endhead
\param{contours} & From \param{Slice}. & required \\
\param{printCenter} & Target point on the bed. & (339, $-193$, 0), the usual spot on the lab bed \\
\param{zOffset} & Extra height after placing. & 0.0 \\
\param{move} & An additional manual shift vector. & none \\
\param{seamMode} & \param{0} Keep, \param{1} Align, \param{2} Stagger, \param{3} Nearest, \param{4} Guide, \param{5} Track-below. & 1 (Align) \\
\bottomrule
\end{longtable}

The seam group (\param{seamAngle}, \param{seamPoint},
\param{seamCurve}) only matters for modes 1 to 5. Outputs:
\param{moved}, \param{seams}, \param{info}. Whether the placed part is
reachable is checked in one place only, \param{Robot Simulation}
(see below), not here.

\subsection{Print Bed: describes the plate once}

Instead of entering the bed size and print centre into several
components, \param{Print Bed} (\textbf{2 Place}) describes the
plate once (width, depth, centre, and default print centre) and emits a
definition string for \param{Place}, \param{Robot
Simulation}, and \param{Live View}. With a bed connected,
\param{Place} uses the print centre of the bed instead of the
built-in default. Explicitly connected inputs still take precedence.

\subsection{G-code Generator: writes the machine instructions}

The G-code Generator is the largest component. The inputs fall into
three groups: speed and shape, flow and pressure (the levers specific
to paste), and optional feature groups that are added only when needed.

\textbf{Speed and shape:} \param{crvs} (from Place), \param{layerHeight},
\param{width}, \param{ramDiameter} (default \param{94}\,mm; a warning
appears outside 5 to 200\,mm, as a sanity check rather than a hard
limit), \param{pressureCoeff} (\param{1.2}, the overall flow
multiplier), \param{printSpeed} (\param{43}), \param{travelSpeed}
(\param{30}), \param{retractSpeed} (\param{1}),
\param{firstLayerSpeedFactor} (\param{0.5}), \param{firmware} (Duet-RRF
or Marlin; selects the G-code dialect), \param{relativeE}
(\param{true}), \param{zOffset} (a nozzle-to-bed calibration shim with
a default of \param{0}; the slicer already applies the layer height, so
leave this at 0 unless the Z zero of the machine needs a correction),
\param{tol}, and \param{savePath}.

\textbf{Machine frame (\param{robot} input):} without a definition, the
output keeps the legacy coordinate frame of the lab SCARA, a mirror
about $y = -x$, so that the saved lab definitions keep printing
correctly. On a Cartesian printer, connect an \param{Robot
Definition} with the profile \param{generic-cartesian},
\param{generic-duet-cartesian}, or \param{none} to obtain the identity
frame. The active frame is always reported in \param{info}
(\texttt{frame:}), so the choice is visible.

\textbf{Flow and pressure (specific to paste):} \param{pressureOffset}
(plunger pull-back, the paste equivalent of retraction),
\param{primeExtra}, \param{coast}, \param{zHop}, \param{hopSpeed},
\param{footerZLift}, \param{dwell}, \param{leadIn}, \param{leadInBoost},
\param{minLayerTime}, \param{retraction} (an on/off dropdown), and
\param{seamTol}.

\textbf{Optional feature groups, added only when needed:}

\begin{itemize}
  \item \textbf{advance} and \textbf{retractaccel}: firmware acceleration and pressure-advance overrides.
  \item \textbf{arcs} (and \textbf{arcwelder}, which requires \param{arcs}): smooth curved output instead of many straight segments.
  \item \textbf{skirt}: a purge ring before the print.
  \item \textbf{floors}: solid base layers.
  \item \textbf{headeroverride}: replaces the generated header.
  \item \textbf{pump}: for a screw pump instead of a ram or plunger; see \S 6.1 of \emph{reference.tex}. The pump group only calculates a recommended pump setting. The E values of the G-code are not changed.
  \item \textbf{varh} (variable layer-height flow): enable this group after adaptive or conical slicing, because the layer height then varies per point.
  \item \textbf{ecomp} (slope-compensated flow): an experimental correction for extra flow on inclined walls. The group is not calibrated yet; leave the group off until an inclined-wall calibration print has been made (see \S 6).
\end{itemize}

Outputs: \param{gcode} (one item per line), \param{preview} and
\param{travels} (curve trees), \param{length} (mm), \param{volume}
(cm$^3$), \param{time} (min), \param{info}, \param{pump} (a summary
when the pump group is on), \param{start} (the first point), and
\param{startGcode} (a jog-to-start block).

\subsection{G-code Restart and Save G-code}

\begin{itemize}
  \item \param{G-code Restart} resumes a print that failed part-way. The G-code is cut at a Z height, a Duet percentage, or a Rhino document-coordinate XYZ point (for example a Live or Console position output). Point wins over percentage, which wins over Z, when more than one is given; the remainder is rebuilt. Outputs are \param{resume} (the restart position) and the \param{printed} and \param{remaining} curve lists.
  \item \param{Save G-code} writes \param{gcode} to a folder with a timestamped file name. The component writes only when the \param{write} input becomes True. Connect a Button rather than a permanent toggle, so that one click produces one file. The write runs in the background to keep the canvas responsive: the component first reports ``writing'' and sets \param{written = true} on the following solve once the file is on disk.
\end{itemize}

\subsection{Robot Definition and Robot Simulation}

The robot is a selection, not an assumption. \param{Robot
Definition} (\textbf{4 Robot}) picks the machine profile and emits a
definition string that every robot-aware component accepts on the
\param{robot} input. The shipped profiles are \param{none} (toolpath
only), \param{generic-cartesian}, \param{generic-duet-cartesian}, and
\param{potterbot-scara-v4} (the lab SCARA). A text panel with a short
alias such as \texttt{"scara"} or \texttt{"none"} works as well.

\begin{itemize}
  \item \param{Robot Simulation} (\textbf{4 Robot}) visualizes the G-code or toolpath with the printed layer and the travel moves in different colours. When the profile supplies kinematics (SCARA), the component poses the robot model and checks reach and joint limits while the simulation is scrubbed. An invalid move sets \param{reachOK} to false, lists the positions in \param{problems}, and raises a red error on the component (``targets are outside robot limits''). Never run a program that shows that error on a real arm. With the profile \param{none}, the component is a plain toolpath and extrusion viewer.
\end{itemize}

Definitions that still use the pre-2.0 component \param{Robot}
continue to open and solve, but that component is no longer on the
ribbon. New work should use Robot Definition and Robot Simulation.

\subsection{Overhang and Print Stats (Analyze)}

\begin{itemize}
  \item \param{Overhang} flags layers and segments whose local angle exceeds a warning threshold (\param{warnDeg}, default 25\iamdeg) or a failure threshold (\param{failDeg}, default 45\iamdeg), colour-coded in the viewport. Both numbers are placeholders and are not calibrated for a specific material; see \S 6.
  \item \param{Print Stats} reports the total length, the extruded volume and mass, and an estimated print time with a per-layer breakdown. The report is a useful check before a long print. With a \param{robot} definition connected, the time estimate also accounts for the joint-speed limits of the machine (a SCARA near the reach boundary is slower than the Cartesian numbers suggest). The robot-limited time is reported next to the Cartesian estimate.
\end{itemize}

\subsection{DoE Design, RSM Evaluate, Desirability, and Test Geometry (Calibrate)}

The calibration components find print settings for a new material by
designing a small set of test prints, measuring the results, and
fitting a model. The reference (\S 6) describes the full workflow and
the exercises contain a guided first pass. In short: \textbf{DoE
Design} creates a run plan, the runs are printed with \textbf{Test
Geometry} coupons and measured by hand, \textbf{RSM Evaluate} fits a
model to the measurements, and \textbf{Desirability} recommends the
best settings.

Measurement tables can be CSV or Excel files. CSV files in English
(\texttt{1.5}) or German (\texttt{1,5}) number format are detected per
file. A file that mixes both conventions is refused with the offending
line named and is never read as wrong numbers.

% =====================================================================
\section{Checking the output before printing}
% =====================================================================

\begin{enumerate}
  \item Read the \param{info} output of every component that was changed. The components report problems per object rather than failing silently.
  \item Inspect \param{preview} (extruding moves) and \param{travels} (non-extruding moves) in the viewport.
  \item Run \param{Print Stats} on the finished G-code and check \param{timeMin}, \param{volumeCm3}, and \param{massG} against the expected values.
  \item On a robot arm, run \param{Robot Simulation} with the matching \param{Robot Definition} connected. \param{reachOK} must be true, \param{problems} must be empty, and no red error may be shown before anything runs on the machine.
  \item Only then save the file (\param{Save G-code} or the \param{savePath} input) and send the file to the machine. Watch the first layer in person.
\end{enumerate}

% =====================================================================
\section{Placeholder values}
% =====================================================================

A few defaults are not calibrated. They exist so that the plugin works
out of the box, not because they are the right numbers for a specific
material:

\begin{itemize}
  \item \param{warnDeg} and \param{failDeg} on \param{Overhang} (25\iamdeg{} and 45\iamdeg).
  \item The pump calibration CSV of the pump group. The shipped file is a placeholder; the bucket-test protocol in \S 6.1 of \emph{reference.tex} produces the real table for a material and pump.
  \item \param{eCompRefMul} of the \param{ecomp} group (slope-compensated flow). Leave the group off until an inclined-wall calibration print has been made.
\end{itemize}

None of these values damage anything when left at the default. They are
not accurate for a specific material until calibrated.

% =====================================================================
\section{Non-planar modes and robot arms}
% =====================================================================

The conical mode of \param{Slice}, and robot simulation, are
advanced topics. \emph{reference.tex} covers the failure modes and
caveats. Two points matter from the start:

\begin{itemize}
  \item Conical mode assumes that the nozzle never tilts: the layers bend, the tool stays vertical. \param{nozzleClear} is the collision guard; raising it hides a real risk instead of removing one.
  \item \param{Robot Simulation} raises a red error for moves outside the limits of the machine, on the SCARA profile. Never override or ignore that error. (The pre-2.0 robot component exposed the same check as a \param{danger} output.)
\end{itemize}

% =====================================================================
\section{Further reading}
% =====================================================================

\begin{itemize}
  \item The exercises (\emph{exercises.tex}) build up the components one at a time in a Grasshopper-only dry run, ending in a shared print project.
  \item The reference (\emph{reference.tex}) serves as a lookup during work and explains the reasoning behind each setting and the edge cases that this guide does not cover.
\end{itemize}

\end{document}
